\documentclass{plantillaPPT}

\titulo{10}{Convergencia de Series}

\begin{document}
\primeraDiapo



\begin{frame}{Criterio de Divergencia}
\framesubtitle{Criterios de Convergencia}

Dada una serie \(\displaystyle\sum_{n}^{\infty}a_n\), sabemos que:\\[2ex]
\[\text{Si} \sum_{n}^{\infty}a_n\ \text{es Convergente} \Rightarrow \lim_{n\to\infty}a_n = 0\]\\[2ex]
Luego, por contrarrecíproco, obtenemos el siguiente criterio de convergencia:
\[\text{Si} \lim_{n\to\infty}a_n \neq 0 \Rightarrow \sum_{n}^{\infty}a_n\ \text{Es Divergente}\]
    
\end{frame}

\begin{frame}{Criterio de Comparación Directa}
\framesubtitle{Criterios de Convergencia}

Dadas las sucesiones \(a_n\) y \(b_n\) tales que \(\forall n\in\mathbb{N}\) \(0\leq a_n\leq b_n\) entonces tiene que:\\[3ex]
\begin{center}
- Si \(\displaystyle\sum_{n}^\infty b_n\) CV \(\Rightarrow\) \(\displaystyle\sum_{n}^\infty a_n\) CV\\[3ex]
- Si \(\displaystyle\sum_{n}^\infty a_n\) DV \(\Rightarrow\) \(\displaystyle\sum_{n}^\infty b_n\) DV
\end{center}
    
\end{frame}

\begin{frame}{Acotamientos y expresiones útiles}
\framesubtitle{Protips}

Las siguientes expresiones son útiles al acotar funciones:
\vspace{5pt}
\begin{itemize}
    \item Desigualdad triangular: \(|a\pm b| \leq |a|+|b|\)
    \vspace{5pt}
    \item Dados \(a,b>0\): \(\frac{1}{a+b}\leq\frac{1}{a}\) o bien \(\frac{1}{a+b}\leq\frac{1}{b}\)
    \vspace{5pt}
    \item $|\sen(x)|\leq1$
    \item $|\cos(x)|\leq1$
\end{itemize}
    
\end{frame}

\begin{frame}{Comparación al límite}
\framesubtitle{Criterios de Convergencia}

Dados \(a_n \geq 0\) y \(b_n>0\) tales que:
\[\lim_{n\to\infty}\frac{a_n}{b_n}=L\]
- Si \(L\in\mathbb{R} \land L\neq0\) entonces la series \(\displaystyle\sum_{n}^{\infty}a_n\) y \(\displaystyle\sum_{n}^{\infty}b_n\) comparten convergencia(ambas divergen o convergen).\\[2ex]
- \(L = 0\): \(\quad\) Si \(\displaystyle\sum_{n}^{\infty}b_n\) CV \(\Rightarrow\) \(\displaystyle\sum_{n}^{\infty}a_n\) CV. \\[2ex]
- \(L = \infty\): \(\quad\) Si \(\displaystyle\sum_{n}^{\infty}b_n\) DV \(\Rightarrow\) \(\displaystyle\sum_{n}^{\infty}a_n\) CV.

\end{frame}

\begin{frame}{Criterio de la Integral}
\framesubtitle{Criterios de Convergencia}

Sea \(a_n\) una sucesión, y \(f\) una función continua, positiva no creciente en el intervalo \([1,\infty[\) que cumple \(f(k) = a_k, \; \forall k\in\mathbb{N}\). Entonces:\\[4ex]
\begin{center}
\(\displaystyle\sum_{n}^{\infty}a_n\) CV \(\Longleftrightarrow\) \(\displaystyle\int_1^{\infty}f(x)\,dx\) CV
\end{center}
\end{frame}

\begin{frame}{Criterio de Leibniz}
\framesubtitle{Criterios de Convergencia}

Sea \(a_n\) una sucesión, si se cumplen las siguientes condiciones:\\[2ex]
- \(|a_n|\) es decreciente.\\[2ex]
- \(\displaystyle\lim_{n\to\infty}|a_n| = 0\)\\[2ex]
Entonces se cumple que las series:\\[2ex]
\[\sum_{n}^\infty (-1)^{n} a_n \quad \text{o bien} \quad \sum_{n}^\infty (-1)^{n+1} a_n\]\\[2ex]
Son convergentes.

\end{frame}

\begin{frame}{Criterio de Convergencia Absoluta}
\framesubtitle{Criterios de Convergencia}

Dada una sucesión \(a_n\) se tiene que:\\[2ex]
Si \(\displaystyle\sum_{n}^\infty |a_n|\) CV \(\Rightarrow\) \(\displaystyle\sum_{n}^\infty a_n\) CV\\[2ex]
Y se dice que la serie \(\displaystyle\sum_{n}^\infty a_n\) es \textbf{absolutamente convergente}.\\[4ex]
Además, si \(\displaystyle\sum_{n}^\infty a_n\) CV pero \(\displaystyle\sum_{n}^\infty |a_n|\) DV, se dice que la serie \(\displaystyle\sum_{n}^\infty a_n\) es \textbf{Condicionalmente convergente}.
\end{frame}

\begin{frame}{Criterio de la Razón/Cociente y de la Raíz}
\framesubtitle{Criterios de Convergencia}

Sea \(\displaystyle\sum_{n}^\infty a_n\) una serie de términos no nulos que cumple alguna de las expresiones:\\[2ex]

\[\lim_{n\to\infty}\frac{|a_{n+1}|}{|a_n|} = L \qquad \text{o bien} \qquad \lim_{n\to\infty}\sqrt[n]{|a_n|} = L\]\\[2ex]
Se tiene que:\\[2ex]
- Si \(L<1\) La serie Converge, y converge \textbf{absolutamente}. \\[2ex]
- Si \(L>1\) La serie Diverge.\\[2ex]
- Si \(L=1\) El criterio no es concluyente.
    
\end{frame}

\begin{frame}{Problema 1}
\framesubtitle{Práctica 10}
1. Considere la suma infinita
\begin{equation*}
    54+18+6+2+\frac{2}{3}+\frac{2}{9}+...
\end{equation*}
Escriba como serie geométrica y determine su suma.
\end{frame}

\begin{frame}{Problema 2}
\framesubtitle{Práctica 10}

2. Determinar la suma de las siguientes series convergente.
\begin{equation*}
    (a) \sum_{n=5}^\infty \frac{3^{n-1}}{4^{3n+1}} \qquad \qquad (b) \sum_{n=2}^\infty \frac{1}{n^2-1}
\end{equation*}
\end{frame}

\begin{frame}{Problema 3}
\framesubtitle{Práctica 10}
3. Analice la convergencia de las siguientes series:
\begin{equation*}
(a) \sum_{n=1}^\infty \arcsen\left(\frac{1}{\sqrt{n}}\right) \qquad \qquad (b) \sum_{n=1}^\infty \frac{3^n}{n^24^n+\sqrt{n}}
\end{equation*}
\end{frame}

\begin{frame}
\framesubtitle{Práctica 10}
\begin{equation*}
    (c) \sum_{n=1}^\infty \left(1+\frac{1}{n}\right)^n \qquad \qquad (d)
\sum_{n=3}^\infty \frac{\ln(n)}{n^{3/2}}
\end{equation*}
\end{frame}

\begin{frame}{Problema 4}
\framesubtitle{Práctica 10}
4. Analice la convergencia de las siguientes series alternantes.
\begin{equation*}
(a)\sum_{n=1}^\infty \frac{\cos(n\pi)(n+2)}{n^3+n^2+1} \qquad \qquad (b) \sum_{n=2}^\infty (-1)^n\,\frac{\sqrt{n^2+1}}{n^3}
\end{equation*}
    
\end{frame}


\end{document}